\subsection{Fixed-target metric realization and its sharp dimensional scope}
\label{wb:sphere-metric-limit-note}

This note concerns the codomains of the three based maps
\[
 q_2^{\wedge 2}:X^{\wedge 2}\cong S^{12}
   \longrightarrow Z_{12}:=\mathcal A^{\wedge 2},
\]
\[
 q_{1,\mathrm{cusp}}^{\wedge 4}:X^{\wedge 4}\cong S^{24}
   \longrightarrow Z_{24}^{\mathcal B}:=\mathcal B^{\wedge 4},
 \qquad
 q_1^{\wedge 4}:X^{\wedge 4}\cong S^{24}
   \longrightarrow Z_{24}^{\mathcal C}:=\mathcal C^{\wedge 4}.
\]
In particular, \(S^{12}\) and \(S^{24}\) are the sources, not the spaces
in which the metric limits occur.  We construct one fixed compatible
length metric on each codomain which realizes the native
four-dimensional flat limits on the target slices.  The realization is
sharp: after the same rescaling, the full ambient limits have dimensions
\(8\) and \(12\), and the collapsed smash base point goes to infinity.

Put
\[
 r_T=e^{-2\pi T},\qquad
 J=\begin{pmatrix}0&1\\-1&0\end{pmatrix},\qquad
 Y_T=TJ+E_T,
\]
where \(E_T\) is the bounded matrix in the radial cusp expansion.  Choose
and fix \(T_0\) and \(C>0\) so that
\(\|E_T\|_{\mathrm{op}}\leq C\) for every \(T\geq T_0\), in the fixed
Euclidean coordinates used below.  This notation avoids confusing the
error matrix with the quotient circle \(C_T\) used below.  On the
fixed torus
\[
 \overline T^2=(\overline\Lambda\otimes\mathbb R)/\overline\Lambda
              \cong\mathbb R^2/\mathbb Z^2
\]
let
\[
 h_T:=\iota_T^*g_{Q_T}=dx^{\mathsf T}Y_T^{\mathsf T}Y_Tdx.
\]
Thus
\begin{equation}
 T^{-2}h_T\longrightarrow h_\infty:=dx^{\mathsf T}dx
 \label{wb:sphere-metric-normalized-Q-limit}
\end{equation}
in every \(C^k\) norm in the fixed flat coordinates.  Indeed, the cusp
coefficients are holomorphic in \(t=e^{-2\pi T}\), so their \(T\)-derivatives
are exponentially small, while \(T^{-1}Y_T=J+T^{-1}E_T\).

Use normalized angular coordinates
\[
 u,\theta\in\mathbb R/\mathbb Z;
\]
\(u\) is the argument of the base coordinate and \(\theta\) is the value of
the invariant character.  Increase this same \(T_0\) once more, if
necessary, so that all radial cusp formulas used below hold for
\(T\geq T_0\) and \(e^{-2\pi T_0}<\epsilon\), and put
\[
 r_0=e^{-2\pi T_0}.
\]
For \(0<r\leq r_0\), write
\[
 T(r):=-\frac{1}{2\pi}\log r.
\]

\paragraph{Target-slice notation.}
We use the embeddings
\(j_T,k_T^{\mathcal B},k_T^{\mathcal C}\) defined in
\cref{wb:cusp-cone-target-slice-embeddings}; the normalized angle \(u=0\) used
below is the same positive radial base angle denoted by \(1\in S^1\) there.

\begin{theorem}[Topology-forced collapse for every fixed compatible metric]
\label{wb:sphere-metric-uniform-collapse}
Let \(D_{12},D_{24}^{\mathcal B},D_{24}^{\mathcal C}\) be arbitrary fixed
metrics compatible with the respective compact target topologies, and let
\(*\) denote the corresponding smash base point.  Then
\begin{align}
 d_H^{D_{12}}\bigl(j_T(Q_T^2),\{*\}\bigr)&\longrightarrow0,\\
 d_H^{D_{24}^{\mathcal B}}
  \bigl(k_T^{\mathcal B}(C_T^4),\{*\}\bigr)&\longrightarrow0,\\
 d_H^{D_{24}^{\mathcal C}}
  \bigl(k_T^{\mathcal C}(C_T^4),\{*\}\bigr)&\longrightarrow0.
 \label{wb:sphere-metric-Hausdorff-collapse}
\end{align}
Consequently, with the metrics restricted from any one fixed ambient
target, all three target slices converge in Gromov--Hausdorff topology to
a point, not to \(\mathbb R^4\).
\end{theorem}

\begin{proof}
We prove the first assertion; the other two are identical.  Let
\[
 G:[0,r_0]^2\times(T^3)^2\longrightarrow Z_{12}
\]
be the product of the two inner-cone quotient maps followed by the smash
quotient.  It is continuous and
\(G(0,0,z_1,z_2)=*\) for every \((z_1,z_2)\in(T^3)^2\).  Since the domain
is compact, uniform continuity with respect to \(D_{12}\) gives
\[
 \sup_{(z_1,z_2)\in(T^3)^2}
 D_{12}\bigl(G(r,r,z_1,z_2),*\bigr)\longrightarrow0.
\]
The set \(j_T(Q_T^2)\) is contained in such a cross-section, which proves
the first Hausdorff convergence.  Replacing \(T^3\) by \(T^2\) and taking
four factors proves the other two.  Hausdorff convergence to a singleton
implies Gromov--Hausdorff convergence to that singleton.
\end{proof}

The preceding theorem is unrescaled.  It does not prevent a blow-up of one
fixed metric.  The following construction gives the strongest natural
positive statement.

\begin{theorem}[One fixed cusp-adapted metric realizes the vertical
\(\mathbb R^4\) limits]
\label{wb:sphere-metric-fixed-realization}
There are fixed compatible length metrics
\[
 D_{12}\text{ on }Z_{12},\qquad
 D_{24}^{\mathcal B}\text{ on }Z_{24}^{\mathcal B},\qquad
 D_{24}^{\mathcal C}\text{ on }Z_{24}^{\mathcal C}
\]
with the following properties.

On the inner cone of \(\mathcal A\), the metric is induced by
\begin{equation}
 g_{\mathcal A}
 =dr^2+r^2du^2+\frac{r^2}{T(r)^2}h_{T(r)}.
 \label{wb:sphere-metric-A-fixed-metric}
\end{equation}
On the cusp cone of either \(\mathcal B\) or \(\mathcal C\), it is induced
by
\begin{equation}
 g_{\mathcal D}=dr^2+r^2(du^2+d\theta^2),
 \qquad\mathcal D\in\{\mathcal B,\mathcal C\}.
 \label{wb:sphere-metric-BC-fixed-metric}
\end{equation}
The metrics on the smash powers are the metric quotients of the product
length metrics after the closed fat wedge is collapsed.  Write
\(g_{Z_{12}}\) and \(g_{Z_{24}^{\mathcal D}}\) for the resulting smooth
Riemannian tensors on the open regular strata.

Let
\[
 p_T^{12}:=j_T(0,0),\qquad
 p_T^{24,\mathcal D}:=k_T^{\mathcal D}(0,0,0,0).
\]
Put
\[
 \Sigma_T^{12}:=j_T(Q_T^2),\qquad
 \Sigma_T^{24,\mathcal D}:=k_T^{\mathcal D}(C_T^4).
\]
Let \(d_{T,12}^{\mathrm{slice}}\) and
\(d_{T,24,\mathcal D}^{\mathrm{slice}}\) be the intrinsic Riemannian path
metrics of the first fundamental forms on these embedded slices.  These
must be distinguished from the restricted ambient distance functions
\[
 d_{T,12}^{\mathrm{amb}}
 :=D_{12}|_{\Sigma_T^{12}\times\Sigma_T^{12}},\qquad
 d_{T,24,\mathcal D}^{\mathrm{amb}}
 :=D_{24}^{\mathcal D}|_{\Sigma_T^{24,\mathcal D}
                         \times\Sigma_T^{24,\mathcal D}}.
\]
In general the latter can be smaller, because an ambient path is allowed
to leave the slice.

Then the induced first fundamental forms on the target slices satisfy the
exact identities
\begin{equation}
 \left(\frac{T}{r_T}\right)^2
  j_T^*g_{Z_{12}}
 =g_{Q_T}\oplus g_{Q_T},
 \label{wb:sphere-metric-A-slice-exact}
\end{equation}
and
\begin{equation}
 \left(\frac{\ell_T}{r_T}\right)^2
  (k_T^{\mathcal D})^*g_{Z_{24}^{\mathcal D}}
 =g_{C_T}^{\oplus4}.
 \label{wb:sphere-metric-BC-slice-exact}
\end{equation}
Consequently,
\begin{align}
 \left(\Sigma_T^{12},\frac{T}{r_T}d_{T,12}^{\mathrm{slice}},
       p_T^{12}\right)
 &\longrightarrow(\mathbb R^4,0),\\
 \left(\Sigma_T^{24,\mathcal D},
       \frac{\ell_T}{r_T}d_{T,24,\mathcal D}^{\mathrm{slice}},
       p_T^{24,\mathcal D}\right)
 &\longrightarrow(\mathbb R^4,0)
 \label{wb:sphere-metric-slice-fixed-blowups}
\end{align}
in pointed local-isometric and pointed Gromov--Hausdorff topology, and
smoothly on the exhausting Euclidean slice balls.

The full ambient pointed blow-ups are instead
\begin{align}
 \left(Z_{12},\frac{T}{r_T}D_{12},p_T^{12}\right)
 &\longrightarrow(\mathbb R^8,0),\\
 \left(Z_{24}^{\mathcal D},
       \frac{\ell_T}{r_T}D_{24}^{\mathcal D},
       p_T^{24,\mathcal D}\right)
 &\longrightarrow(\mathbb R^{12},0).
 \label{wb:sphere-metric-ambient-fixed-blowups}
\end{align}
as pointed metric spaces in pointed Gromov--Hausdorff topology.  More
precisely, for every fixed radius the corresponding rescaled target ball is
eventually contained in an open regular-stratum coordinate neighborhood,
and the pulled-back Riemannian tensors converge there in every \(C^k\) norm.
Thus smooth convergence is asserted on an exhaustion by regular-stratum
neighborhoods, not globally on the singular compact targets.  Under these
limits, the slice \(\mathbb R^4\) is respectively the coordinate subspace
obtained by fixing the two radial and two base-angle coordinates, or by
fixing the four radial and four base-angle coordinates.
\end{theorem}

\begin{proof}
Equation \eqref{wb:sphere-metric-normalized-Q-limit} shows that the metrics on
the compact cross-sections of \eqref{wb:sphere-metric-A-fixed-metric} remain
uniformly bounded and converge to a positive definite flat metric.
Therefore the cross-section diameter is \(O(r)\).  Moreover the radial
coordinate is \(1\)-Lipschitz and a radial segment from level \(r\) to the
vertex has length \(r\).  The length completion consequently adds exactly
one point at \(r=0\), and its topology is the prescribed cone topology.
The same assertions are immediate for
\eqref{wb:sphere-metric-BC-fixed-metric}.  Interpolation on a closed annulus
away from the vertex extends these metrics to fixed compatible length
metrics on the compact spaces \(\mathcal A,\mathcal B,\mathcal C\).

If a compact metric space \((E,d)\) and a closed subset \(W\subset E\) are
given, the metric quotient which collapses \(W\) to one point is
\[
 d_{E/W}([x],[y])
 =\min\{d(x,y),d(x,W)+d(y,W)\}.
\]
Applying this to the fat wedges in the product length spaces gives the
claimed fixed compatible metrics on the smash targets.

On the \(r_T\)-slice of \(\mathcal A\), with \(u=0\), the first fundamental
form is
\[
 \frac{r_T^2}{T^2}h_T.
\]
The definition \(h_T=\iota_T^*g_{Q_T}\) proves
\eqref{wb:sphere-metric-A-slice-exact}.  On the \(r_T\)-slice of
\(\mathcal B\) or \(\mathcal C\), with \(u=0\), the form is
\(r_T^2d\theta^2\).  Under
\(\vartheta_T:[\theta]\mapsto[\ell_T\theta]\), the native quotient-circle
metric is \(\ell_T^2d\theta^2\).  This proves
\eqref{wb:sphere-metric-BC-slice-exact}.  The native systole estimates
\[
 \operatorname{sys}(Q_T^2)\geq T-C,
 \qquad
 \operatorname{sys}(C_T^4)=\ell_T
\]
then prove \eqref{wb:sphere-metric-slice-fixed-blowups}: every prescribed
radius ball is eventually exactly a Euclidean four-ball.

It remains to compute the ambient limits.  In one \(\mathcal A\)-factor,
put
\[
 \lambda_T=\frac{T}{r_T},\qquad
 s=\lambda_T(r-r_T),\qquad
 \eta=Tu,\qquad z=Y_Tx
\]
on the universal cover of the torus directions.  On every fixed bounded
set in \((s,\eta,z)\), one has
\[
 \frac r{r_T}=1+\frac sT,
 \qquad
 T(r)=T+O(T^{-1}),
 \qquad
 Y_{T(r)}Y_T^{-1}=I+O(T^{-2}),
\]
with the same estimates after any fixed number of derivatives.  Thus
\(\lambda_T^2g_{\mathcal A}\) converges in these coordinates to
\[
 ds^2+d\eta^2+dz_1^2+dz_2^2.
\]
The inward radial boundary is at distance \(T\), the \(u\)-period is \(T\),
and the \(z\)-period lattice is \(Y_T\mathbb Z^2\), whose systole is at
least \(T-C\).  The outward edge of the fixed cusp chart is also at scaled
distance
\[
 \frac{T}{r_T}(r_0-r_T)
 =T\left(\frac{r_0}{r_T}-1\right)\longrightarrow\infty.
\]
Thus the inward vertex, all periods, and the outward edge recede to
infinity.  A single factor converges on exhausting regular-stratum
coordinate neighborhoods to \(\mathbb R^4\), and two factors converge to
\(\mathbb R^8\).

For one \(\mathcal B\)- or \(\mathcal C\)-factor, put
\[
 \lambda_T=\frac{\ell_T}{r_T},\qquad
 s=\lambda_T(r-r_T),\qquad
 \eta=\ell_Tu,\qquad \zeta=\ell_T\theta.
\]
Then
\[
 \lambda_T^2g_{\mathcal D}
 =ds^2+\left(\frac r{r_T}\right)^2
       (d\eta^2+d\zeta^2),
\]
which converges smoothly on bounded coordinate sets to the Euclidean metric
on \(\mathbb R^3\).  The two circle periods and the inward radial boundary
are all at scaled distance \(\ell_T\to\infty\), while the outward edge is
at scaled distance
\[
 \frac{\ell_T}{r_T}(r_0-r_T)
 =\ell_T\left(\frac{r_0}{r_T}-1\right)\longrightarrow\infty.
\]
Four factors therefore converge on exhausting regular-stratum coordinate
neighborhoods to \(\mathbb R^{12}\).

Finally, a level-\(r_T\) factor point has distance exactly \(r_T\) from its
cone vertex.  Indeed, the radial segment has that length.  A competing
curve contained in the inner chart has length at least \(r_T\), while one
which leaves has inner-chart portions of total length at least
\((r_0-r_T)+r_0>r_T\).  The Riemannian-\(\ell^2\)-product distance from a
synchronized point to the fat wedge is the minimum of the factor
distances, and hence is also \(r_T\).

For the chosen scale put
\[
 b_T=\lambda_Tr_T,\qquad a_T=\lambda_T(r_0-r_T),
\]
and let \(\sigma_T\) be the least scaled period in a factor.  Thus
\(\sigma_T\geq T-C\) in the \(\mathcal A\) calculation and
\(\sigma_T=\ell_T\) in the \(\mathcal B,\mathcal C\) calculation.  The
preceding computations give
\[
 a_T\longrightarrow\infty,\qquad
 b_T\longrightarrow\infty,\qquad
 \sigma_T\longrightarrow\infty.
\]
Fix \(R\).  A curve of rescaled length at most \(R\) from \(p_T\) meets
neither the fat wedge nor \(r=r_0\) for all large \(T\), so the
radius-\(R\) target ball lies in the open regular stratum.  If \(x,y\) lie
in that ball, a competitor through the fat wedge has length at least
\[
 (b_T-R)+(b_T-R)>2R,
\]
and a competitor which exits the inner chart and returns has length at
least
\[
 2(a_T-R)>2R.
\]
Every remaining competitor of length at most \(2R\) stays in
\(r/r_T=1+o_R(1)\).  After lifting it to the angular and torus covers, a
nonzero deck translation has length at least
\((1-o_R(1))(\sigma_T-2R)>2R\) for all large \(T\).  Since the path through
\(p_T\) has length at most \(2R\), none of these three kinds of global
shortcut computes the distance between \(x\) and \(y\).  Thus the metric
on every fixed target ball is computed entirely in one lifted
regular-stratum chart.  The \(C^\infty\) convergence of the pulled-back
tensors on these exhausting charts implies the pointed
Gromov--Hausdorff limits in
\eqref{wb:sphere-metric-ambient-fixed-blowups}.
\end{proof}

\begin{theorem}[Complete intrinsic scalar-rescaling trichotomy]
\label{wb:sphere-metric-scaling-trichotomy}
For the fixed metrics of
\cref{wb:sphere-metric-fixed-realization}, let \(\lambda_T>0\) be arbitrary
and put
\[
 \beta_T:=\lambda_Tr_T.
\]
Choose one of the following triples, retaining the indicated ambient
metric:
\[
 \begin{split}
  (\Sigma_T,d_T^{\mathrm{slice}},d_T^{\mathrm{amb}},p_T;D)
  &=(\Sigma_T^{12},d_{T,12}^{\mathrm{slice}},
      d_{T,12}^{\mathrm{amb}},p_T^{12};D_{12}),\\
  (\Sigma_T,d_T^{\mathrm{slice}},d_T^{\mathrm{amb}},p_T;D)
  &=(\Sigma_T^{24,\mathcal D},
      d_{T,24,\mathcal D}^{\mathrm{slice}},
      d_{T,24,\mathcal D}^{\mathrm{amb}},
      p_T^{24,\mathcal D};D_{24}^{\mathcal D}),
      \qquad \mathcal D\in\{\mathcal B,\mathcal C\}.
 \end{split}
\]
After passing to a subsequence so that \(\beta_T\) has a limit in
\([0,\infty]\), the intrinsic pointed slice limit is
\begin{equation}
 (\Sigma_T,\lambda_Td_T^{\mathrm{slice}},p_T)
 \longrightarrow
 \begin{cases}
  \{\mathrm{pt}\},&\beta_T\longrightarrow0,\\
  ((\mathbb R/L\mathbb Z)^4,0),
      &\beta_T\longrightarrow L\in(0,\infty),\\
  (\mathbb R^4,0),&\beta_T\longrightarrow\infty.
 \end{cases}
 \label{wb:sphere-metric-scaling-trichotomy-result}
\end{equation}
The convergence in
\eqref{wb:sphere-metric-scaling-trichotomy-result} is pointed
Gromov--Hausdorff in all three cases.  In the finite positive case the flat
metrics converge smoothly under the fixed marked-torus identifications; in
the unrolling case the convergence is pointed local-isometric and smooth on
the exhausting Euclidean balls.
In the middle line the lattice is orthogonally equivalent to the square
lattice \(L\mathbb Z^4\).  At the same time,
\begin{equation}
 \lambda_TD(p_T,*)=\beta_T
 \label{wb:sphere-metric-vertex-distance}
\end{equation}
for every one of the displayed choices.

Only in the unrolling regime \(\beta_T\to\infty\) do we compare the
intrinsic metric with the restricted ambient distance.  In that regime, for
every \(R\in(0,\infty)\),
\begin{equation}
 \sup_{x,y\in
  \overline B_{\lambda_Td_T^{\mathrm{slice}}}(p_T,R)}
 \left|
  \lambda_Td_T^{\mathrm{amb}}(x,y)
  -\lambda_Td_T^{\mathrm{slice}}(x,y)
 \right|\longrightarrow0.
 \label{wb:sphere-metric-intrinsic-ambient-agreement}
\end{equation}
Moreover, there are \(\varepsilon_{T,R}\to0\) such that
\begin{equation}
 \overline B_{\lambda_Td_T^{\mathrm{slice}}}(p_T,R)
 \subset
 \overline B_{\lambda_Td_T^{\mathrm{amb}}}(p_T,R)
 \subset
 \overline B_{\lambda_Td_T^{\mathrm{slice}}}
             (p_T,(1+\varepsilon_{T,R})R).
 \label{wb:sphere-metric-intrinsic-ambient-ball-comparison}
\end{equation}
Thus the restricted ambient metric has the same pointed
\(\mathbb R^4\) limit on these fixed slice balls.  No agreement is asserted
in the regimes \(\beta_T\to0\) or \(\beta_T\to L\in(0,\infty)\), where
radial or fat-wedge shortcuts need not disappear.  Every rescaling which
unrolls the intrinsic slice sends the cone/fat-wedge base point to infinity,
so no such limit is a tangent cone based at that point.

In the same regime \(\beta_T\to\infty\), the corresponding full ambient
bubbles are
\begin{align}
 (Z_{12},\lambda_TD_{12},p_T^{12})
 &\longrightarrow(\mathbb R^8,0),\\
 (Z_{24}^{\mathcal D},\lambda_TD_{24}^{\mathcal D},
  p_T^{24,\mathcal D})
 &\longrightarrow(\mathbb R^{12},0),
 \qquad\mathcal D\in\{\mathcal B,\mathcal C\}.
 \label{wb:sphere-metric-arbitrary-unrolling-ambient-bubbles}
\end{align}
These are pointed Gromov--Hausdorff limits of the singular metric targets.
The convergence of tensors is \(C^\infty\) only under coordinate
identifications on exhausting open regular-stratum neighborhoods; no
global smooth convergence of the singular targets is claimed.
\end{theorem}

\begin{proof}
For the \(Q_T^2\) route, the rescaled period matrix in each two-dimensional
factor is
\[
 \frac{\lambda_Tr_T}{T}Y_T
 =\beta_T\left(J+\frac{E_T}{T}\right).
\]
If \(\beta_T\to L\in(0,\infty)\), this lattice converges to
\(LJ\mathbb Z^2\); two factors give the square four-torus in
\eqref{wb:sphere-metric-scaling-trichotomy-result}.  If
\(\beta_T\to0\), the lengths of a lattice basis tend to zero, so the
diameter tends to zero.  If \(\beta_T\to\infty\), the reverse triangle
inequality gives systole
\[
 \frac{\beta_T}{T}(T-C)
 =\beta_T\left(1-\frac CT\right)\longrightarrow\infty,
\]
so every fixed-radius ball is eventually Euclidean.

For a \(C_T^4\) slice, the four periods after rescaling have the common
length
\[
 \lambda_Tr_T=\beta_T,
\]
which gives the same three intrinsic cases directly.

Consider first one factor.  A curve from a level-\(r_T\) point to the cone
vertex which remains in the inner chart has length at least \(r_T\), because
the \(dr^2\)-term makes \(r\) \(1\)-Lipschitz there.  If it leaves the
chart, it has a first crossing of \(r=r_0\) and, before ending at the
vertex, a last re-entry through \(r=r_0\); its two inner-chart portions
have total length at least
\((r_0-r_T)+r_0>r_T\).  The radial segment has length \(r_T\), so the
factor distance to its based vertex is exactly \(r_T\).  The
Riemannian-\(\ell^2\)-product distance from every synchronized slice point
to the fat wedge is the minimum of these factor distances and hence is
also \(r_T\).  The metric-quotient formula in the preceding proof therefore
gives \eqref{wb:sphere-metric-vertex-distance}.

It remains to prove the restricted-ambient comparison, so assume
\(\beta_T\to\infty\) and fix \(R\).  Two points in
\(\overline B_{\lambda_Td_T^{\mathrm{slice}}}(p_T,R)\) have intrinsic
distance at most \(2R/\lambda_T\).  The slice path is an ambient competitor,
so
\[
 d_T^{\mathrm{amb}}(x,y)\leq d_T^{\mathrm{slice}}(x,y).
\]
The quotient shortcut through the fat wedge has rescaled length
\(2\lambda_Tr_T=2\beta_T\), and hence is longer than \(2R\) for all large
\(T\).  A path which leaves the fixed cusp chart must reach \(r=r_0\).
Its outward cost is controlled by
\begin{equation}
 \lambda_T(r_0-r_T)
 =\beta_T\left(\frac{r_0}{r_T}-1\right)\longrightarrow\infty,
 \label{wb:sphere-metric-outward-shortcut-exclusion}
\end{equation}
so no path of rescaled length at most \(2R\) can use the interpolating
annulus or the global part of a factor.

Consider now an ambient competitor of length at most
\(2R/\lambda_T+o(\lambda_T^{-1})\) which stays in the cusp chart.  In every
factor its radial coordinate satisfies
\[
 |r-r_T|\leq \frac{2R+o(1)}{\lambda_T},
 \qquad
 \frac r{r_T}=1+O_R(\beta_T^{-1}).
\]
On this tube, the horizontal tensors in the directions tangent to the slice
are uniformly \((1+\varepsilon_{T,R})\)-bilipschitz to their values at
\(r=r_T\), where \(\varepsilon_{T,R}\to0\).  For
\(\mathcal B\) and \(\mathcal C\) this follows directly from
\((r/r_T)^2=1+O_R(\beta_T^{-1})\).  For \(\mathcal A\), the same conclusion
follows, without changing the fixed coordinates or the matrix \(Y_T\), from
\[
 \begin{aligned}
  T(r)-T&=O_R(\beta_T^{-1}),\\
  \frac{T(r)}T&=1+O_R((\beta_TT)^{-1}),\\
  Y_{T(r)}Y_T^{-1}&=I+O_R((\beta_TT)^{-1}).
 \end{aligned}
\]
Projecting the competitor to the torus or character directions, and
discarding its radial and base-angle components, produces a path between
the same slice endpoints.  The bilipschitz estimate therefore gives
\[
 d_T^{\mathrm{slice}}(x,y)
 \leq(1+\varepsilon_{T,R})d_T^{\mathrm{amb}}(x,y).
\]
Together with the opposite inequality and
\(\lambda_Td_T^{\mathrm{slice}}(x,y)\leq2R\), this proves
\eqref{wb:sphere-metric-intrinsic-ambient-agreement}.  The first inclusion in
\eqref{wb:sphere-metric-intrinsic-ambient-ball-comparison} follows from
\(d_T^{\mathrm{amb}}\leq d_T^{\mathrm{slice}}\).  For a point \(x\) in
the middle ball, apply the same fat-wedge and outward exclusions to a
near-minimizing ambient path from \(p_T\) to \(x\), whose rescaled length is
at most \(R+o(1)\), and project that path to the slice directions.  This
gives
\[
 \lambda_Td_T^{\mathrm{slice}}(p_T,x)
 \leq(1+\varepsilon_{T,R})R
\]
after enlarging \(\varepsilon_{T,R}\) by a quantity tending to zero, and
proves the second inclusion.

For completeness, the full ambient bubbles at an arbitrary unrolling scale
are obtained without changing any coordinate or orientation.  In one
\(\mathcal A\)-factor use
\[
 s=\lambda_T(r-r_T),\qquad
 \eta=\beta_Tu,\qquad
 z=\frac{\beta_T}{T}Y_Tx.
\]
Writing \(A_T(r)=Y_{T(r)}Y_T^{-1}\), one obtains the exact expression
\[
 \lambda_T^2g_{\mathcal A}
 =ds^2+\left(\frac r{r_T}\right)^2d\eta^2
  +\left(\frac r{r_T}\right)^2
   \left(\frac{T}{T(r)}\right)^2
   dz^{\mathsf T}A_T(r)^{\mathsf T}A_T(r)dz.
\]
On every fixed bounded coordinate set,
\[
 \frac r{r_T}=1+O_R(\beta_T^{-1}),\qquad
 T(r)-T=O_R(\beta_T^{-1}),\qquad
 A_T(r)=I+O_R((\beta_TT)^{-1}),
\]
with the same estimates after every fixed number of derivatives.  Hence
the tensor converges smoothly to the Euclidean metric on \(\mathbb R^4\).
The \(\eta\)-period is \(\beta_T\), and the \(z\)-period lattice is
\[
 \frac{\beta_T}{T}Y_T\mathbb Z^2,
 \qquad
 \operatorname{sys}\left(
   \frac{\beta_T}{T}Y_T\mathbb Z^2\right)
 \geq\beta_T\left(1-\frac CT\right)\longrightarrow\infty.
\]
In one \(\mathcal B\)- or \(\mathcal C\)-factor use
\[
 s=\lambda_T(r-r_T),\qquad
 \eta=\beta_Tu,\qquad \zeta=\beta_T\theta.
\]
Then
\[
 \lambda_T^2g_{\mathcal D}
 =ds^2+\left(\frac r{r_T}\right)^2
       (d\eta^2+d\zeta^2),
\]
and both angular periods are \(\beta_T\to\infty\).

In all three cases the inward boundary and the fat wedge are at scaled
distance \(\beta_T\), while a lift starting at radius \(r_T\) incurs the
one-way radial cost in
\eqref{wb:sphere-metric-outward-shortcut-exclusion} before its first outward
crossing of \(r=r_0\).  Fix \(R\).  For large \(T\), all these quantities
and all nonzero period lengths exceed \(4R\).
If \(x,y\) lie in the ambient rescaled radius-\(R\) ball, their distances
from the fat wedge are at least \(\beta_T-R\); hence a quotient path through
the smash base point has length at least \(2(\beta_T-R)>2R\).  Every
remaining competitor lifts to the complement of the fat wedge in the
product.  If such a lift exits the inner product chart, then some factor
has a first crossing of \(r=r_0\); the radial variation before that crossing
is at least \(\lambda_T(r_0-r_T)-R\), and returning to the lift of \(y\)
incurs the same lower bound.  Its length is therefore greater than \(2R\).
Finally, a path of length at most \(2R\) remains in the tube
\(r/r_T=1+O_R(\beta_T^{-1})\); there every nonzero deck translation has
length tending to infinity, so such a path cannot wrap around an angular
or torus period.  The displayed coordinate neighborhoods therefore exhaust
the pointed balls without quotient or global shortcuts.
Their smooth tensor convergence proves the pointed Gromov--Hausdorff
limits in
\eqref{wb:sphere-metric-arbitrary-unrolling-ambient-bubbles}, with smooth
convergence confined to this regular-stratum exhaustion.
\end{proof}

\begin{corollary}[Unit-volume normalization does not give
\(\mathbb R^4\)]
\label{wb:sphere-metric-unit-volume}
The four-volume of \(Q_T^2\) is \((\det Y_T)^2\sim T^4\).  Its unit-volume
length factor is therefore \((\det Y_T)^{-1/2}\sim T^{-1}\), and its
unit-volume limit is a square flat four-torus, not \(\mathbb R^4\).
Likewise, the four-volume of \(C_T^4\) is \(\ell_T^4\), and its unit-volume
normalization is exactly the square unit four-torus.
\end{corollary}

\begin{proposition}[What vertex tangents and compact asymptotic cones can
mean]
\label{wb:sphere-metric-tangent-scope}
The compact target topology alone determines no metric tangent cone.
For an exact conic length metric built from fixed piecewise-Riemannian
metrics on the links, however, the vertex links have topological and
Hausdorff dimensions \(7\) for \(Z_{12}\) and \(11\) for both
\(Z_{24}\) targets.  The resulting vertex cones have topological and
Hausdorff dimensions \(8\) and \(12\), and are not \(\mathbb R^4\).

More precisely, writing \(L_3=T^3\) and \(L_2=T^2\), one has the based
topological models
\[
 Z_{12}\cong\operatorname{Susp}(L_3*L_3),
 \qquad
 Z_{24}^{\mathcal B}
 \cong\operatorname{Susp}(L_2*L_2*L_2*L_2).
\]
At the suspension base point their conic links are respectively
\(L_3*L_3\) and \(L_2^{*4}\).  Also
\[
 (\mathcal C,c_0)
 \cong
 (\mathcal M/\partial\mathcal M,[\partial\mathcal M]),
 \qquad
 (Z_{24}^{\mathcal C},*)
 \cong
 (\mathcal M^4/\partial(\mathcal M^4),
  [\partial(\mathcal M^4)]),
\]
so an exact conic metric at the collapsed boundary has link
\(\partial(\mathcal M^4)\).  At the collapsed point of
\(Z_{24}^{\mathcal C}\), the complete integral local-homology calculation is
\begin{equation}
 H_i\bigl(Z_{24}^{\mathcal C},
          Z_{24}^{\mathcal C}\setminus\{*\};\mathbb Z\bigr)
 \cong
 \begin{cases}
  \mathbb Z^{\binom4{i-1}},&2\leq i\leq5,\\
  0,&i=6,7,\\
  \mathbb Z^{\binom4{i-8}},&8\leq i\leq12,\\
  0,&\text{otherwise}.
 \end{cases}
 \label{wb:sphere-metric-C-smash-local-homology}
\end{equation}
In particular, its local homology is not concentrated in top degree.  For
the first two targets the top local homology occurs in degree \(8\) and
\(12\), respectively; for the third target
\eqref{wb:sphere-metric-C-smash-local-homology} includes a top
\(\mathbb Z\) in degree \(12\).  None has the local homology of
\(\mathbb R^4\).

Finally, the standard asymptotic cone obtained by blowing down any fixed
compact metric target is a point.  Indeed, if \(R_n\to\infty\), then
\[
 \operatorname{diam}(Z,R_n^{-1}D)
 =R_n^{-1}\operatorname{diam}(Z,D)\longrightarrow0.
\]
Thus the limits in \cref{wb:sphere-metric-fixed-realization} are
moving-base-point blow-ups, often called bubbles.  They are neither
vertex-based tangent cones nor asymptotic cones of the compact targets.
\end{proposition}

\begin{proof}
For compact spaces \(V_i\) with closed subspaces \(A_i\), there is a
canonical homeomorphism
\[
 (V_1/A_1)\wedge\cdots\wedge(V_m/A_m)
 \cong
 \frac{V_1\times\cdots\times V_m}
      {\bigcup_iV_1\times\cdots\times A_i\times\cdots\times V_m}.
\]
Take \(V_i=\operatorname{Cone}(L_i)\) and \(A_i=L_i\), its outer boundary.
The product of \(m\) cones is homeomorphic to the cone on the iterated join
\(L_1*\cdots*L_m\), and its displayed union of faces is the outer link.
Collapsing that link gives its suspension.  This proves the first two
models and the link dimensions
\[
 \dim(T^3*T^3)=3+3+1=7,
 \qquad
 \dim((T^2)^{*4})=4\cdot2+3=11.
\]

We next prove the based homeomorphism for \(\mathcal C\), rather than
identifying the two presentations without a map.  Put
\(L=\partial\mathcal M\), write the attached cone as
\[
 \operatorname{Cone}(L)
 =L\times[0,1]/(L\times\{0\}),
\]
with base \(L\times\{1\}\), and choose a collar
\(\kappa:L\times[0,1]\hookrightarrow\mathcal M\) satisfying
\(\kappa(x,0)=x\).  There is a homeomorphism
\[
 \Phi:\mathcal C\setminus\{c_0\}
       \longrightarrow\operatorname{int}\mathcal M
\]
given on the punctured cone, the inner collar, and the remaining core by
\[
 \Phi([x,t])=\kappa(x,t/e)\quad(0<t\leq1),\qquad
 \Phi(\kappa(x,s))=\kappa(x,e^{s-1})\quad(0\leq s\leq1),
\]
\[
 \Phi(q)=q
 \quad\text{for}\quad
 q\in\mathcal M\setminus\kappa(L\times[0,1)).
\]
The first two formulas agree at the glued points
\([x,1]=\kappa(x,0)\), and the second and third agree at
\(\kappa(x,1)\).  Their collar-coordinate images are respectively
\(L\times(0,e^{-1}]\) and \(L\times[e^{-1},1]\); hence the displayed
piecewise map and its piecewise inverse are continuous.

The compact Hausdorff spaces \(\mathcal C\) and
\(\mathcal M/L\) are both one-point compactifications of
\(\operatorname{int}\mathcal M\): for \(\mathcal C\) use \(\Phi\), and for
\(\mathcal M/L\) use the quotient map on
\(\mathcal M\setminus L=\operatorname{int}\mathcal M\).  Uniqueness of the
one-point compactification therefore extends \(\Phi\) to a based
homeomorphism
\begin{equation}
 (\mathcal C,c_0)
 \xrightarrow{\ \cong\ }
 (\mathcal M/\partial\mathcal M,[\partial\mathcal M]).
 \label{wb:sphere-metric-C-boundary-quotient-homeomorphism}
\end{equation}
Only after establishing
\eqref{wb:sphere-metric-C-boundary-quotient-homeomorphism} do we apply the
quotient identity at the start of the proof.  It gives
\[
 \mathcal C^{\wedge4}
 \cong
 \frac{\mathcal M^4}
 {\displaystyle\bigcup_{a=1}^4
  \mathcal M^{a-1}\times\partial\mathcal M
  \times\mathcal M^{4-a}}
 =\mathcal M^4/\partial(\mathcal M^4),
\]
and carries the smash base point to the collapsed boundary.

Put \(N=\mathcal M^4\).  Rounding its product corners preserves the
homeomorphism type of the pair \((N,\partial N)\) and makes
\(\partial N\) an eleven-dimensional closed manifold.  A boundary collar
descends in \(N/\partial N\) to a cone neighborhood of the collapsed
point, so the local link is exactly \(\partial N\) and
\begin{equation}
 H_i(N/\partial N,(N/\partial N)\setminus\{*\};\mathbb Z)
 \cong\widetilde H_{i-1}(\partial N;\mathbb Z).
 \label{wb:sphere-metric-boundary-link-local-shift}
\end{equation}

The integral homology calculation in the proof of
\cref{wb:cusp-cone-S24-route} for the specific compact orientable
three-manifold \(\mathcal M\) gives
\[
 H_j(\mathcal M;\mathbb Z)\cong
 \begin{cases}
  \mathbb Z,&j=0,1,\\
  0,&\text{otherwise}.
 \end{cases}
\]
All groups are free, so the integral Kunneth theorem gives
\[
 H_j(N;\mathbb Z)\cong
 \begin{cases}
  \mathbb Z^{\binom4j},&0\leq j\leq4,\\
  0,&\text{otherwise}.
 \end{cases}
\]
Poincare--Lefschetz duality, with the product orientation retained, gives
\[
 H_j(N,\partial N;\mathbb Z)
 \cong H^{12-j}(N;\mathbb Z)
 \cong
 \begin{cases}
  \mathbb Z^{\binom4{j-8}},&8\leq j\leq12,\\
  0,&\text{otherwise}.
 \end{cases}
\]
The long exact sequence of \((N,\partial N)\) now has disjoint absolute
and relative degree ranges.  It gives
\[
 \widetilde H_j(\partial N;\mathbb Z)\cong
 \begin{cases}
  \mathbb Z^{\binom4j},&1\leq j\leq4,\\
  0,&j=5,6,\\
  \mathbb Z^{\binom4{j-7}},&7\leq j\leq11,\\
  0,&\text{otherwise}.
 \end{cases}
\]
Substitution in
\eqref{wb:sphere-metric-boundary-link-local-shift} proves
\eqref{wb:sphere-metric-C-smash-local-homology}.

An exact metric cone is invariant under every radial blow-up at its vertex,
and its Hausdorff dimension is one plus that of its fixed
piecewise-Riemannian link.  For the first target,
\[
 H_8\bigl(\operatorname{Cone}(T^3*T^3),T^3*T^3;\mathbb Z\bigr)
 \cong\mathbb Z;
\]
the analogous top local group for the \(\mathcal B\) target is
\(\mathbb Z\) in degree \(12\), while the full \(\mathcal C\) calculation
is \eqref{wb:sphere-metric-C-smash-local-homology}.  Euclidean
\(\mathbb R^4\) has local homology \(\mathbb Z\) only in degree \(4\), so
none of these conic vertex models is \(\mathbb R^4\).  The diameter
calculation proves the final assertion.
\end{proof}

\paragraph{Consequences for the three sphere routes.}
The exact positive statement is that one fixed metric on each compact
codomain contains the native \(\mathbb R^4\) decompactification as a
vertical four-plane in a moving-point Euclidean bubble.  The whole bubbles
are \(\mathbb R^8\) for the \(S^{12}\) route and \(\mathbb R^{12}\) for the
\(S^{24}\) routes.  The factor \(\beta_T=\lambda_Tr_T\) proves that this
dimension discrepancy cannot be removed by a different scalar rescaling
of the same fixed cusp-adapted metric: the slice unrolls only when the
smash base point has already gone to infinite scaled distance.

No metric on a source sphere is supplied by these quotient constructions.
If an independent smooth Riemannian metric is placed on \(S^{12}\) or
\(S^{24}\), its ordinary tangent spaces are \(\mathbb R^{12}\) or
\(\mathbb R^{24}\).  Pulling a codomain metric back through
\(q_2^{\wedge2}\) or \(q_1^{\wedge4}\), on the regular strata where the
maps are differentiable, gives a degenerate quadratic form along the
positive-dimensional quotient kernels, not a Riemannian source metric.

The moving four-sphere carriers
\[
 \mathbb S_{12,T}^4=S_T^2\wedge S_T^2\subset S^6,\qquad
 \mathbb S_{24,T}^4=C_T^{\wedge4}\subset S^8
\]
of \cref{wb:cusp-cone-angle-forgotten-sphere-carriers} carry separately
assigned metrics.  Namely, \(S_T^2\) first receives the quotient length
metric obtained from
\((\overline T^2,h_T)\cong(Q_T,g_{Q_T})\) by collapsing
\(K^{(1)}\subset\overline T^2\), equivalently by collapsing
\(\iota_T(K^{(1)})\subset Q_T\).  Then
\(\mathbb S_{12,T}^4\) receives the quotient length metric of the
Riemannian \(\ell^2\)-product \((S_T^2)^2\) after its fat wedge is
collapsed.  Likewise, \(\mathbb S_{24,T}^4\) receives the quotient length
metric of the Riemannian \(\ell^2\)-product \(C_T^4\) after its fat wedge is
collapsed.  These metrics are assigned directly to the moving carriers;
they are unrescaled intrinsic quotient metrics, not restrictions of
\(D_{12}\) or \(D_{24}^{\mathcal B}\), whose domains are the conical
targets, and no fixed ambient metric on \(S^6\) or \(S^8\) is used to
define them.  Their Euclidean balls are
based at the midpoint classes
\[
 \begin{aligned}
  y_*&=[(1/2,1/2)]\in\overline T^2,\qquad z_T=c(y_*),\\
  \zeta_{12,T}&=[z_T,z_T],\\
  \zeta_{24,T}
  &=[(\ell_T/2,\ell_T/2,\ell_T/2,\ell_T/2)].
 \end{aligned}
\]
not at the images of \(p_T^{12}\) or \(p_T^{24,\mathcal B}\), which belong
to the subsequently collapsed loci.  In particular, the trichotomy in
\cref{wb:sphere-metric-scaling-trichotomy} is the trichotomy of the
uncollapsed intrinsic flat-torus slices; it is not a finite-\(\beta_T\)
description of the singular quotient-metric moving four-spheres.  Indeed,
the restriction of any one compatible metric on the compact sphere
\(S^6\) or \(S^8\) has diameter bounded by the diameter of that sphere.  By
contrast, in the separately assigned moving metrics the midpoint distances
to the collapsed loci are at least \((T-C)/2\) and exactly
\(\ell_T/2\), respectively, and hence tend to infinity.  Thus the assigned
unrescaled metrics cannot themselves be restrictions of one fixed ambient
metric.  This does not preclude a separately proved realization after a
scalar blow-up of another fixed quotient metric.

Finally, the rank obstruction for a full-base torus quotient must retain
its actual metric hypothesis: \emph{quotient systole tends to infinity}
means the quotient norm induced by the unrescaled Euclidean fibre metric.
Without that qualification the assertion is false, since an arbitrary
\(T\)-dependent scalar multiplication can make the systole of any fixed
nonzero quotient tend to infinity.

\subsection{Expanded proof of the retained-angle cusp class}
\label{wb:h24:ret:section}

\paragraph{The exact object.}
We use the fixed based smooth identification of the family member $X$ with
$S^6$, its cusp filling, and the complete original toric coordinates of
\cref{mapsI:src:A.1.2}. The native quotient and compact extension are
\cref{wb:cusp-cone-native-descent,wb:cusp-cone-S12-route}.
The local parameter $t$ in the following formulas is the cusp parameter
$q$ of the family convention; $B_0=J$ in the radial metric notation.
All original period terms and native lattices are retained. The
calculation is independent of the angle-forgetting map to $S^3$.

\paragraph{Original coordinates and the specified sphere-valued map}

The ordered quotient lattice is
\[
 \overline\Lambda=\mathbb Z\langle\overline{\widehat\gamma},
              \overline{\widehat u}\rangle,
 \qquad
 K_\Lambda=\mathbb Z\langle\widehat w,\widehat\delta\rangle,
 \qquad
 \overline T^2=\mathbb R^2/\mathbb Z^2.
\]
In these ordered bases the original period data are
\begin{equation}
 B_0=\begin{pmatrix}0&1\\-1&0\end{pmatrix},\qquad
 Z(s)=sB_0+C(t),\qquad t=e^{2\pi i s},\qquad
 C(t)=\begin{pmatrix}6\mu(t)&h_{\rm per}(t)\\
 b(t)-h_{\rm per}(t)&\mu(t)\end{pmatrix}.
 \label{wb:h24:ret:eq:original-periods}
\end{equation}
Here \(b=\beta+\tau\), \(s=\tau-h_{\rm per}(t)\), and every entry of
\(C\) is holomorphic at \(t=0\), exactly as in the original source.
The subscript in \(h_{\rm per}\) distinguishes that period function
from the sphere-valued map below.  Write
\begin{equation}
 R(t)=-2\pi\operatorname{Im}C(t),\qquad
 \ell=\log|t|,\qquad A(t)=\ell B_0+R(t).
 \label{wb:h24:ret:eq:A}
\end{equation}
Choose a closed cusp disc \(|t|\leq\epsilon\) inside the original
coordinate disc and small enough for the quotient construction.  On a
slightly larger closed disc put \(M=\sup\|R(t)\|<\infty\).  We will
choose a regular value with \(0<r_0<\epsilon\) and
\(|\log r_0|>M\).  This chooses a point of the already specified map;
it does not change that map or the cutoff radius \(\epsilon\).

The original covering model on its dense torus is
\begin{equation}
 x_j=e^{2\pi i\zeta_j},\qquad
 \Psi_n(x,t)=
 \bigl(e^{2\pi i C(t)n}t^{B_0n}x,t\bigr),\qquad n\in\mathbb Z^2.
 \label{wb:h24:ret:eq:deck}
\end{equation}
For a branch of \(s\), writing
\(\zeta=Z(s)u+k\) with \(u,k\in\mathbb R^2\) gives
\begin{equation}
 \log|x|=A(t)u.
 \label{wb:h24:ret:eq:logx}
\end{equation}
Indeed, \(\operatorname{Im}s=-\ell/(2\pi)\), so
\(-2\pi\operatorname{Im}(Z(s)u+k)=\ell B_0u+R(t)u\).
Since \(B_0\) is orthogonal,
\begin{equation}
 \bigl|\bigl(\ell B_0+vR(t)\bigr)\xi\bigr|
 \geq (|\ell|-vM)|\xi|>0
 \quad(0\leq v\leq1,\ \xi\ne0)
 \label{wb:h24:ret:eq:invertibility}
\end{equation}
whenever \(|\ell|>M\).  Thus \(A(t)\) is invertible there.  Equation
\eqref{wb:h24:ret:eq:deck} changes \(A(t)^{-1}\log|x|\) by \(n\), so the native
quotient is exactly
\begin{equation}
 \overline\pi_f([x,t])=[A(t)^{-1}\log|x|]\in\mathbb R^2/\mathbb Z^2.
 \label{wb:h24:ret:eq:native}
\end{equation}
This also follows from the flat marking in \cref{wb:cusp-cone-native-descent}.

Let
\[
 T^3=S^1_\phi\times\overline T^2,
 \qquad S^1_\phi=\mathbb R/(2\pi\mathbb Z),
\]
with coordinate basepoints \(\phi=0\), \(u_1=0\), \(u_2=0\).
Its coordinate two-skeleton is the following specified closed subset:
\begin{equation}
 K^{(2)}=
 \bigl(\{0\}\times\overline T^2\bigr)
 \cup\bigl(S^1_\phi\times\{u_1=0\}\bigr)
 \cup\bigl(S^1_\phi\times\{u_2=0\}\bigr).
 \label{wb:h24:ret:eq:two-skeleton}
\end{equation}
The product CW structure has one open three-cell, with coordinates
\((\phi,u_1,u_2)\in(0,2\pi)\times(0,1)^2\).  Consequently
\[
 c_3:T^3\longrightarrow T^3/K^{(2)}\cong S^3
\]
is the degree-one quotient, where the orientation of \(S^3\) is fixed
by the displayed coordinate order.  One explicit identification of its
open cell with \(\mathbb R^3\) is
\[
 (\phi,u_1,u_2)\longmapsto
 \bigl(-\cot(\phi/2),-\cot(\pi u_1),-\cot(\pi u_2)\bigr).
\]
Each derivative in its own variable is strictly positive, and approach
to any boundary face leaves every compact subset of \(\mathbb R^3\).
Extending by its point at infinity gives the stated based sphere
identification and verifies its continuity on the collapsed skeleton.

The exact conical target and map are
\begin{align}
 \mathcal A&=
 \frac{[0,\epsilon]\times T^3}
 {\{0\}\times T^3\sim v_0,\quad
  \{\epsilon\}\times T^3\sim v_1},
 \label{wb:h24:ret:eq:A-target}\\
 q_2(x)&=[r,e^{i\phi},\overline\pi_f(x)]
       &&(t(f(x))=re^{i\phi},\ 0<r<\epsilon),\notag\\
 q_2(W)&=v_0,\qquad
 q_2\bigl(X\setminus\operatorname{int}N_0\bigr)=v_1.
 \label{wb:h24:ret:eq:q2}
\end{align}
Fix the same source basepoint \(x_{\rm base}\in W\).  Define
\begin{equation}
 p_3:\mathcal A\longrightarrow\Sigma_r(S^3,*),\qquad
 p_3([r,z])=[r,c_3(z)],\qquad
 H=p_3q_2:(X,x_{\rm base})\longrightarrow(S^4,*).
 \label{wb:h24:ret:eq:H}
\end{equation}
The reduced suspension here collapses both endpoint copies of \(S^3\)
and the entire interval over \(*\) to one point.  Thus \(p_3\) sends
both \(v_0\) and \(v_1\) to \(*\) and is well-defined and based.
It is continuous by the defining quotient topology.  An explicit
identification of its open four-cell with \(\mathbb R^4\) is
\begin{equation}
 (r,\phi,u_1,u_2)\longmapsto
 \bigl(-\cot(\pi r/\epsilon),-\cot(\phi/2),
       -\cot(\pi u_1),-\cot(\pi u_2)\bigr).
 \label{wb:h24:ret:eq:four-cell}
\end{equation}
This fixes the target orientation.  Replacing \(r\) locally by
\(\ell=\log r\) preserves this orientation, since \(dr/d\ell=r>0\).
The ordered coordinates in \eqref{wb:h24:ret:eq:four-cell} specify the sphere
identification used throughout this calculation.

\paragraph{The statement proved.} We give the complete expanded proof of \cref{maps:src:B.3.1}, including the full differential, the relative smoothing and the spin bordism calculation.


\paragraph{The entire regular inverse image lies in one original chart}

Choose the actual lift
\begin{equation}
 u_0=(-1/4,1/4)^{\mathsf T},\qquad
 p_0=B_0u_0=(1/4,1/4)^{\mathsf T},\qquad \phi_0=\pi.
 \label{wb:h24:ret:eq:chosen-value}
\end{equation}
Modulo \(\mathbb Z^2\), \(u_0\) is \((3/4,1/4)\) in the open cell.  Put
\(t_0=r_0e^{i\phi_0}\) and
\(y_0=[r_0,c_3(e^{i\phi_0},[u_0])]\).  The numbers \(r_0\) and
\(t_0\) will remain nonzero throughout the framing calculation.

The original triangle \(\{0,e_1,e_2\}\) has toric chart
\begin{equation}
 \mathcal U_\sigma\cong\mathbb C^3,\qquad
 (z_0,z_1,z_2)=(t/(x_1x_2),x_1,x_2),\qquad
 t=z_0z_1z_2.
 \label{wb:h24:ret:eq:z-chart}
\end{equation}
These characters are dual to \((0,0,1),(1,0,1),(0,1,1)\), respectively,
as direct pairing verifies.  This is the chart of
\cref{mapsI:src:A.1.2}, with its original coordinates.

For \(j=1,2\), introduce the nowhere vanishing holomorphic functions
\begin{equation}
 h_j(t)=\exp\bigl(-2\pi i(C(t)u_0)_j\bigr).
 \label{wb:h24:ret:eq:units}
\end{equation}
The following is an exact coordinate map, with the exact inverse
displayed alongside it:
\begin{align}
 F(z)&=(w_0,w_1,w_2)
 =\bigl(z_0/(h_1(t)h_2(t)),\ h_1(t)z_1,\ h_2(t)z_2\bigr),
 \label{wb:h24:ret:eq:F}\\
 F^{-1}(w)&=
 \bigl(w_0h_1(t)h_2(t),\ w_1/h_1(t),\ w_2/h_2(t)\bigr),
 \quad t=w_0w_1w_2.
 \label{wb:h24:ret:eq:Finv}
\end{align}
Both maps are holomorphic wherever \(|t|\) is in the period disc,
both preserve \(t\), and substitution proves that both composites are
the identity.  In particular the maps extend across the origin and
across every coordinate divisor.  No period matrix is replaced by a
different matrix in this coordinate change.

\begin{lemma}\label{wb:h24:ret:chart-lemma}
For sufficiently small \(r_0>0\), the whole inverse image
\(M=H^{-1}(y_0)\) is represented injectively inside one contractible
coordinate neighborhood in \(X\).  In its \(w\)-coordinates it is
exactly the embedded two-torus
\begin{equation}
 w(a,b)=
 \bigl(\rho_0e^{i(\phi_0-a-b)},\rho_1e^{ia},\rho_2e^{ib}\bigr),
 \qquad
 \rho_0=r_0^{1/2},\quad \rho_1=\rho_2=r_0^{1/4},
 \quad (a,b)\in\bigl(\mathbb R/(2\pi\mathbb Z)\bigr)^2.
 \label{wb:h24:ret:eq:M}
\end{equation}
\end{lemma}
\begin{proof}
The point \(y_0\) lies off the collapsed subset of the target, so
\eqref{wb:h24:ret:eq:H} and \eqref{wb:h24:ret:eq:q2} imply
\[
 M=\{x:t(f(x))=t_0,\ \overline\pi_f(x)=[u_0]\}.
\]
Any lift of such a point to \((\mathbb C^*)^2\times\{t_0\}\) has
\(A(t_0)^{-1}\log|x|=u_0+n\) for one \(n\in\mathbb Z^2\).
Applying the actual transformation \(\Psi_{-n}\) from
\eqref{wb:h24:ret:eq:deck} produces a lift with
\(\log|x|=A(t_0)u_0\).  If two such lifts represent the same point,
\(\Psi_n\) relates them and changes their lifted quotient coordinate
by \(n\); hence \(n=0\), and the lifts are identical.  This proves
both completeness and injectivity of this chosen representative.

Since
\(\log|h_j(t)|=2\pi\operatorname{Im}(C(t)u_0)_j
=-(R(t)u_0)_j\), its \(w\)-coordinates satisfy
\[
 (\log|w_1|,\log|w_2|)^{\mathsf T}
 =A(t_0)u_0-R(t_0)u_0=\ell_0p_0,
 \qquad \ell_0=\log r_0.
\]
Thus \(|w_1|=|w_2|=r_0^{1/4}\), and the unchanged product
\(w_0w_1w_2=t_0\) forces \(|w_0|=r_0^{1/2}\) and the displayed
phase.  Conversely every pair \((a,b)\) in \eqref{wb:h24:ret:eq:M} gives one
such lift.  The original kernel coordinates \(k_j\) had period one
in \(x_j=e^{2\pi i\zeta_j}\); the phase coordinates \(a,b\) have
period \(2\pi\), with no cover or change of lattice.

The covering \(Y_\epsilon\to N_0\) is injective on a neighborhood of
the origin in \eqref{wb:h24:ret:eq:z-chart}, because the original action is free
and properly discontinuous.  Choose a ball \(V\) centered at that
origin with closure inside such a neighborhood and the period disc.
Choose a ball \(V_w\) about the origin with closure contained in
\(F(V)\).  Every coordinate in \eqref{wb:h24:ret:eq:M} tends uniformly to zero
as \(r_0\downarrow0\).  Therefore for sufficiently small \(r_0\)
the entire torus lies in \(V_w\).  The inverse image of \(V_w\)
under \(F\), followed by the injective covering projection, is the
required contractible chart in \(X\).
\end{proof}

\paragraph{The full differential and its ordered normal framing}

We calculate in a neighborhood of the compact torus \eqref{wb:h24:ret:eq:M}.
The quotient coordinates admit a unique local lift near \(u_0\).
Set
\[
 v=(\log|w_1|,\log|w_2|)^{\mathsf T},\quad
 \ell=\log|w_0w_1w_2|,\quad
 \phi=\arg(w_0w_1w_2)
\]
using the branch of \(\phi\) near \(\phi_0=\pi\).
Equations \eqref{wb:h24:ret:eq:logx} and \eqref{wb:h24:ret:eq:F} give the exact identity
\begin{equation}
 v=\ell B_0u+R(t)(u-u_0),\qquad
 u=u_0+A(t)^{-1}(v-\ell p_0).
 \label{wb:h24:ret:eq:exact-u}
\end{equation}
Differentiate this identity before restricting to the inverse image:
\begin{equation}
 du=-A(t)^{-1}(dA(t))A(t)^{-1}(v-\ell p_0)
        +A(t)^{-1}(dv-p_0\,d\ell).
 \label{wb:h24:ret:eq:full-derivative}
\end{equation}
On \(M\) the vector \(v-\ell p_0\) is exactly zero.  Therefore
\begin{equation}
 du|_M=A_0^{-1}(dv-p_0\,d\ell),\qquad A_0=A(t_0).
 \label{wb:h24:ret:eq:derivative-on-M}
\end{equation}
This explicitly accounts for all derivatives of the original bounded
period terms; none are omitted or estimated away.

Identify \(T_w\mathbb C^3\) with \(\mathbb C^3\) in the given coordinate chart,
and define real tangent vectors
\begin{equation}
 L_j=(0,\ldots,w_j,\ldots,0),\qquad
 Q_j=(0,\ldots,iw_j,\ldots,0)\qquad(0\leq j\leq2).
 \label{wb:h24:ret:eq:LQ}
\end{equation}
Thus \(d\log|w_i|(L_j)=\delta_{ij}\),
\(d\arg w_i(Q_j)=\delta_{ij}\), and the crossed evaluations vanish.
The ordered tangent vectors of the parametrization \eqref{wb:h24:ret:eq:M} are
\begin{equation}
 T_a=-Q_0+Q_1,\qquad T_b=-Q_0+Q_2.
 \label{wb:h24:ret:eq:tangents}
\end{equation}
They are independent because their \(w_1\) and \(w_2\) components,
respectively, are nonzero.  For the target order
\((\ell,\phi,u_1,u_2)\), define the following four lifts of its
coordinate vectors:
\begin{align}
 N_\ell&=\tfrac12L_0+\tfrac14L_1+\tfrac14L_2,
       &N_\phi&=Q_0,\label{wb:h24:ret:eq:normal-first}\\
 N_i&=-\bigl((A_0)_{1i}+(A_0)_{2i}\bigr)L_0
           +(A_0)_{1i}L_1+(A_0)_{2i}L_2
       &&(i=1,2).\label{wb:h24:ret:eq:normal-u}
\end{align}
Let \(P=(\ell,\phi,u_1,u_2)\).  Evaluation using
\eqref{wb:h24:ret:eq:derivative-on-M} gives the entire differential table
\begin{equation}
\begin{array}{c|rrrrrr}
 &N_\ell&N_\phi&N_1&N_2&T_a&T_b\\\hline
 d\ell&1&0&0&0&0&0\\
 d\phi&0&1&0&0&0&0\\
 du_1&0&0&1&0&0&0\\
 du_2&0&0&0&1&0&0
\end{array}.
 \label{wb:h24:ret:eq:differential-table}
\end{equation}
In particular \(P\), and therefore \(H\) in its open target cell,
is a submersion along \(M\).  The classes of
\((N_\ell,N_\phi,N_1,N_2)\) in \(TX|_M/TM\) give the specified
normal four-framing.  Orthogonal representatives are not required for
this quotient-bundle definition; projecting the displayed lifts along
\(TM\) to any chosen orthogonal complement preserves the quotient
frame.

If the normal frame is specified using the literal four real coordinates
in \eqref{wb:h24:ret:eq:four-cell}, their differential at the chosen value is
\[
 dH_{\rm cell}=K\,dP,\qquad
 K=\operatorname{diag}\!\left(
 \frac{\pi r_0}{\epsilon}\csc^2\frac{\pi r_0}{\epsilon},
 \frac12,2\pi,2\pi\right).
\]
The first factor includes \(dr=r_0\,d\ell\), the second uses
\(\phi_0=\pi\), and the final two use the actual representatives
\((3/4,1/4)\) of \(u_0\). Thus, writing
\(\mathcal N=(N_\ell,N_\phi,N_1,N_2)\), the literal dual frame is
\(\mathcal N K^{-1}\). All four entries of \(K\) are positive
constants on \(M\). The explicitly positive diagonal matrices
\((1-v)K^{-1}+vI_4\), \(0\leq v\leq1\), give a homotopy
from that literal frame to \(\mathcal N\). This retains the original
target coordinates and establishes that either coordinate description
gives the same framed bordism class.

For completeness, the original period matrix admits the following
explicit homotopy of this normal framing.  Put
\begin{equation}
 A_{0,v}=\ell_0B_0+vR(t_0),\qquad 0\leq v\leq1,
 \label{wb:h24:ret:eq:A-homotopy}
\end{equation}
and replace \(A_0\) by \(A_{0,v}\) only in
\eqref{wb:h24:ret:eq:normal-u}.  Equation \eqref{wb:h24:ret:eq:invertibility} proves
invertibility for every \(v\).  Their images under the actual
\((du_1,du_2)\) are the columns of \(A_0^{-1}A_{0,v}\), so these
are normal frames throughout the homotopy.  This is a homotopy of
the computed frame on the same embedded torus, not a replacement of
the original degeneration by a different one.

\begin{lemma}[The exact orientation sign]\label{wb:h24:ret:orientation-lemma}
Use the complex real-coordinate orientation, whose ordered coordinates are
\[
 (\operatorname{Re}w_0,\operatorname{Im}w_0,
  \operatorname{Re}w_1,\operatorname{Im}w_1,
  \operatorname{Re}w_2,\operatorname{Im}w_2).
\]
Then one has
\begin{equation}
 \det(N_\ell,N_\phi,N_1,N_2,T_a,T_b)
       =-\det(A_0)|w_0w_1w_2|^2<0.
 \label{wb:h24:ret:eq:sign}
\end{equation}
Hence the tangent orientation induced by the ordered normal frame is
the orientation of \((-T_a,T_b)\).
\end{lemma}
\begin{proof}
The ordered frame \((L_0,Q_0,L_1,Q_1,L_2,Q_2)\) has determinant
\(|w_0|^2|w_1|^2|w_2|^2>0\) in real coordinates.  Regroup its rows
as \((L_0,L_1,L_2,Q_0,Q_1,Q_2)\); this permutation has three
inversions and sign \(-1\).  Regroup the columns of the frame in
\eqref{wb:h24:ret:eq:sign} as
\((N_\ell,N_1,N_2,N_\phi,T_a,T_b)\); this has two inversions and
sign \(+1\).  The angular block is
\(\left(\begin{smallmatrix}1&-1&-1\\0&1&0\\0&0&1\end{smallmatrix}\right)\),
of determinant one.  The radial block has first column
\((1/2,1/4,1/4)^{\mathsf T}\) and the other columns given by
\eqref{wb:h24:ret:eq:normal-u}.  Adding its last two rows to its first row gives
first row \((1,0,0)\), with remaining lower-right block \(A_0\).
Its determinant is \(\det A_0\).  This proves the equality.
Equation \eqref{wb:h24:ret:eq:A-homotopy} joins \(A_0\) to \(\ell_0B_0\)
through invertible matrices, so
\(\operatorname{sign}\det A_0=\operatorname{sign}\det(\ell_0B_0)=+1\).
The inequality and the tangent orientation follow.
\end{proof}

\paragraph{An explicit contraction of the full frame}

Define a complex-linear matrix, also regarded as a real-linear matrix,
\begin{equation}
 D(a,b)=\operatorname{diag}(e^{-i(a+b)},e^{ia},e^{ib})\in\operatorname{SU}(3).
 \label{wb:h24:ret:eq:D}
\end{equation}
By \eqref{wb:h24:ret:eq:M}, every \(L_j\) and \(Q_j\) at \((a,b)\) is
\(D(a,b)\) times its value at \((0,0)\).  The coefficients in
\eqref{wb:h24:ret:eq:normal-first}, \eqref{wb:h24:ret:eq:normal-u}, and
\eqref{wb:h24:ret:eq:tangents} are constant on \(M\), since \(t=t_0\).
It follows that the positively oriented full frame is exactly
\begin{equation}
 (N_\ell,N_\phi,N_1,N_2,-T_a,T_b)_{(a,b)}
       =D(a,b)G_0,
 \label{wb:h24:ret:eq:full-frame}
\end{equation}
where \(G_0\in\operatorname{GL}^+(6,\mathbb R)\) is its single value at \((0,0)\).
Thus neither an unspecified phase nor a position-dependent
real-linear matrix remains.

Here is a periodic contraction of \(D\), with all parameters
specified.  For \(v\in[0,1]\) put
\[
 s_v=\sin(\pi v/2),\qquad c_v=\cos(\pi v/2),\qquad
 Q_v(\theta)=
 \begin{pmatrix}s_ve^{-i\theta}&-c_v\\c_v&s_ve^{i\theta}\end{pmatrix},
 \qquad K_v(\theta)=Q_v(\theta)Q_v(0)^*.
\]
The two columns of \(Q_v\) have norm one and zero Hermitian inner
product, and \(\det Q_v=s_v^2+c_v^2=1\).  Thus
\(Q_v,K_v\in\operatorname{SU}(2)\).  Direct substitution gives
\[
 K_0(\theta)=I,\quad
 K_1(\theta)=\operatorname{diag}(e^{-i\theta},e^{i\theta}),\quad
 K_v(0)=I.
\]
Let \(\iota_{01},\iota_{02}:\operatorname{SU}(2)\to\operatorname{SU}(3)\) act in the
ordered coordinate planes \((0,1)\) and \((0,2)\), respectively,
and as the identity in the unused coordinate.  Then
\begin{equation}
 D_v(a,b)=\iota_{01}(K_v(a))\,\iota_{02}(K_v(b))
 \label{wb:h24:ret:eq:D-homotopy}
\end{equation}
is a well-defined continuous map on
\([0,1]\times(\mathbb R/(2\pi\mathbb Z))^2\), with
\(D_0=I\) and \(D_1=D\).  This proves the needed null-homotopy
directly, including both phase seams.  In particular
\(D_v(a,b)G_0\) contracts the comparison matrix in
\eqref{wb:h24:ret:eq:full-frame} to one constant positive matrix.

\begin{proposition}\label{wb:h24:ret:framing-proposition}
The stable tangential framing induced by the actual normal four-frame
of \(H^{-1}(y_0)\) agrees, up to homotopy and the orientation sign
\eqref{wb:h24:ret:eq:sign}, with the Lie framing
\((\partial_a,\partial_b)\) on the native two-torus.
\end{proposition}
\begin{proof}
In the coordinate chart, the stable tangent trivialization is the
isomorphism
\[
 \mathbb R^4\oplus TM\longrightarrow T\mathbb C^3|_M,
 \qquad (v,V)\longmapsto
 v_1N_\ell+v_2N_\phi+v_3N_1+v_4N_2+V.
\]
Using the intrinsic frame \((-\partial_a,\partial_b)\) on \(TM\),
its matrix is precisely \eqref{wb:h24:ret:eq:full-frame}.  The homotopy
\eqref{wb:h24:ret:eq:D-homotopy} therefore makes this comparison constant.
A constant positive matrix changes no homotopy class of stable
framing: Gram--Schmidt deforms it to an element of \(SO(6)\), and
an orthogonal matrix of determinant one can be reduced to the identity
by finitely many oriented plane rotations, each joined to the
identity by varying its angle.  This proves the asserted comparison
within the chart.

It remains to compare the chart with the ambient sphere framing.
Pull back the usual stable tangent trivialization
\(TS^6\oplus\mathbb R\cong\mathbb R^7\) under the fixed based diffeomorphism
\(X\cong S^6\).  On the contractible \(w\)-ball containing \(M\),
it differs from the coordinate trivialization plus one real line by
a continuous map into \(\operatorname{GL}(7,\mathbb R)\).  Contracting the ball to its
center contracts this comparison map to a constant.  Thus its
restriction to \(M\) contributes no varying stable framing class.
If the chosen sphere orientation disagrees with the complex
orientation, the remaining constant has negative determinant and
reverses the resulting framed bordism class.  Both that sign and the
explicit sign \eqref{wb:h24:ret:eq:sign} disappear in the order-two class
computed below.  This argument compares actual ambient
trivializations on a contractible chart; it does not suppose that
\(\mathbb C^3\) globally trivializes \(TX\).
\end{proof}

\paragraph{Nonzero framed bordism and the exact unstable class}

We give the nonvanishing test for the resulting torus explicitly.
A stable framing gives a spin structure by lifting its oriented frame
to the trivial \(\operatorname{Spin}\)-bundle.  On the Lie-framed
circle this is the disconnected double cover, whereas the spin
structure induced from the boundary of a framed disc is the
connected double cover: as the boundary tangent turns once, its
\(2\pi\)-rotation lifts from \(1\) to \(-1\) in
\(\operatorname{Spin}(2)\).  We call the former nonbounding and the
latter bounding.

For a closed oriented spin surface \(S\), the spin quadratic
refinement
\[
 q_S:H_1(S;\mathbb F_2)\longrightarrow\mathbb F_2
\]
takes a class represented by an embedded circle to zero or one
according to whether the induced spin structure, using its normal
line in the surface, is bounding or nonbounding.  It satisfies
\begin{equation}
 q_S(x+y)=q_S(x)+q_S(y)+(x\mathbin{\cdot}y).
 \label{wb:h24:ret:eq:quadratic}
\end{equation}
This is the usual spin quadratic form: resolving transverse
intersections adds one tangent turn per intersection, and the
\(2\pi\)-spin-lift calculation above gives the last term.  The
well-defined correspondence between spin structures and such
quadratic refinements is given in
\cite[Theorem 3.2 and Remark 3.3]{DebrayGunningham2018Root}; its original
surface proof is \cite{Johnson1980Root}.

For a symplectic basis \((e_i,f_i)_{i=1}^g\), put
\begin{equation}
 \operatorname{Arf}(S)=\sum_{i=1}^gq_S(e_i)q_S(f_i)\in\mathbb F_2.
 \label{wb:h24:ret:eq:Arf}
\end{equation}
The basis independence can be seen without selecting a preferred
basis: factor the finite sum over the symplectic planes to obtain
\[
 \sum_{x\in H_1(S;\mathbb F_2)}(-1)^{q_S(x)}
   =\prod_{i=1}^g
     \sum_{a,b\in\mathbb F_2}(-1)^{a q_S(e_i)+b q_S(f_i)+ab}
   =2^g(-1)^{\sum_i q_S(e_i)q_S(f_i)}.
\]
The left side has no choice of basis, proving the assertion.

Here is also the bordism-invariance argument needed for the present
obstruction.  A spin bordism has a handle decomposition.  An index
zero or three handle creates or removes a spin sphere, whose first
homology and Arf invariant are zero.  An index one handle either
connects two components, in which case the first-homology quadratic
forms add orthogonally, or increases the genus of one component.
In the latter case the new symplectic pair contains the meridian
of the handle.  That meridian bounds the cocore disc with its
induced spin structure, so its quadratic value is zero; the pair
contributes zero to \eqref{wb:h24:ret:eq:Arf}, regardless of the value on its
dual.  Existing cycles can be represented outside the attaching
discs, so their quadratic values are preserved.  An index two
handle is the reverse of an index one handle.  Therefore each
handle preserves the total Arf invariant of the boundary, and a
spin boundary has Arf invariant zero.  This also proves that a
stably framed boundary has Arf invariant zero, because forgetting
the framing gives a spin bordism.  The full spin bordism statement
is recorded in \cite[Theorem 3.6]{DebrayGunningham2018Root}.

The intrinsic Lie-torus calculation is the known computation in \cite[Example~3.8, p.~6]{DebrayGunningham2018Root}; we retain its direct derivation. For the intrinsic Lie framing on
\(M=(\mathbb R/(2\pi\mathbb Z))^2\), take
\[
 e=[a\mapsto(a,0)],\qquad f=[b\mapsto(0,b)].
\]
Both the tangent and normal vectors in the surface are constant in
the Lie frame along either circle.  Thus both circles have the
nonbounding spin structure, and
\begin{equation}
 e\mathbin{\cdot}f=1,\qquad q_M(e)=q_M(f)=1,
 \qquad q_M(xe+yf)=x+xy+y,
 \qquad \operatorname{Arf}(M)=1.
 \label{wb:h24:ret:eq:actual-Arf}
\end{equation}
Reversing the first Lie tangent vector changes the orientation but
does not change these mod-two values.  In particular the
specifically framed inverse image computed above is not a framed
boundary.

To identify the element, let \(\eta\in\pi_1^{\rm S}\) denote the
stable class of the complex Hopf map.  Its Pontryagin--Thom inverse
image is the circle with its Lie framing.  This can be verified
directly on the Hopf map \(S^3\to\mathbb{CP}^1\): near the fibre
\((z_1,z_2)=(e^{ia},0)\), the target coordinate is \(z_2/z_1\), so
the two horizontal normal vectors are \((0,e^{ia})\) and
\((0,ie^{ia})\).  Add the tangent vector \((ie^{ia},0)\) and the
outward radial vector \((e^{ia},0)\).  In \(\mathbb C^2\) the resulting
full frame makes one complete rotation in each of its two real
coordinate planes.  This loop is contractible explicitly: for the
fixed real-linear involution \(J(z_1,z_2)=(z_1,\overline z_2)\),
\[
 J\operatorname{diag}(e^{ia},e^{ia})J=\operatorname{diag}(e^{ia},e^{-ia})=K_1(-a),
\]
and \(JK_v(-a)J\) contracts the original loop to the identity using
the periodic homotopy already constructed.  Thus the induced stable
tangent frame on the circle is the Lie frame.  This is the basic
framed-circle identification in \cite{Pontryagin1955Root}.

The product of two normally framed inverse images represents the
product in stable homotopy: its tubular collapse is the smash of
the two tubular collapse maps, as follows immediately from the
product tubular neighborhood \(\nu_1\times\nu_2\) and its ordered
normal coordinates.  Consequently the product of the two Lie
circles represents \(\eta^2\).  Equation \eqref{wb:h24:ret:eq:actual-Arf}
proves \(\eta^2\ne0\) by an explicit bordism invariant.  The
stable group \(\pi_2^{\rm S}=\mathbb Z/2\) is Pontryagin's computation
\cite{Pontryagin1950Root,Pontryagin1955Root}; the stable table in
\cite[pp.~384--385]{Hatcher2002} gives the same group.  Thus the sign from
\eqref{wb:h24:ret:eq:sign} satisfies \(-\eta^2=\eta^2\).

\begin{proof}[Expanded proof of \cref{maps:src:B.3.1}]
The globally continuous quotient \(H\) is smooth on a neighborhood
of its compact inverse image \(M\), by
\eqref{wb:h24:ret:eq:exact-u} and \eqref{wb:h24:ret:eq:four-cell}. First make it constant
near the source basepoint, before choosing any approximation bound.
Choose a target chart \(\kappa:O\to\mathbb R^4\) about \(*\) with
\(\kappa(*)=0\), convex ball image, and \(y_0\notin O\).
Choose a source neighborhood \(U\) of \(x_{\rm base}\), disjoint
from \(M\), with \(H(U)\subset O\). Let \(\beta\) be a smooth
function supported in \(U\), equal to one near \(x_{\rm base}\),
with \(0\leq\beta\leq1\). On \(U\) use the based homotopy
\[
 H_s(x)=\kappa^{-1}\bigl((1-s\beta(x))\kappa(H(x))\bigr),
 \qquad 0\leq s\leq1,
\]
and leave \(H\) unchanged off \(U\). Its endpoint \(H_b\) is
smooth and constant near the basepoint; it agrees with \(H\)
near \(M\) and has no new inverse image of \(y_0\).

Choose one closed tubular neighborhood \(U_1\) of \(M\), with
\(M\subset\operatorname{int}U_1\), contained in the unchanged
smooth region and disjoint from the basepoint neighborhood.
Compactness gives the positive bound
\[
 d_0=\min_{x\in X\setminus\operatorname{int}U_1}
             d\bigl(H_b(x),y_0\bigr)>0.
\]
Relative smooth approximation, fixing the same closed set \(U_1\)
and a neighborhood of the basepoint, produces a smooth based map
\(H_{\rm sm}\) whose error from \(H_b\) is less than \(d_0/2\)
on its complement. Choose the approximation closer also than a
fixed target injectivity radius; the short target geodesics give a
based homotopy. Thus the exact inverse image and its differential
are unchanged on \(U_1\), and the complement acquires no new
inverse image. The relative Pontryagin construction applies as in
\cite[\S7, Theorem B and Lemma 4, pp.~43, 48--50]{CuspNullMilnor1965}.

Pontryagin--Thom therefore identifies the stabilization of
\([H]\) with the bordism class of \(M\) carrying the four normal
vectors \eqref{wb:h24:ret:eq:normal-first}--\eqref{wb:h24:ret:eq:normal-u}.  The comparison
\eqref{wb:h24:ret:eq:full-frame} and its explicit contraction identify that
class with the Lie torus, with the harmless order-two orientation
sign already computed.  Hence
\[
 E^\infty[H]=\eta^2\ne0.
\]
Freudenthal suspension is an isomorphism
\[
 E:\pi_{n+2}(S^n)\longrightarrow\pi_{n+3}(S^{n+1})
 \qquad(n\geq4),
\]
because \(n+2\leq2n-2\).  Therefore stabilization from
\(\pi_6(S^4)\) is an isomorphism onto \(\pi_2^{\rm S}\).
The composite \(\eta_4\circ\eta_5:S^6\to S^4\) stabilizes to
\(\eta^2\), so injectivity gives the exact equality
\cref{maps:src:B.3.1}.  Finally, a based null-homotopy of \(q_2\)
would postcompose with \(p_3\) to a based null-homotopy of \(H\),
contradicting \cref{maps:src:B.3.1}.
\end{proof}

\paragraph{Exact scope of the comparison with the angle-forgetting map.}
The map here retains the angular coordinate before collapsing the
coordinate two-skeleton.  Its regular inverse image is the specific
normally four-framed \(T^2\) above.  The angle-forgetting
postcomposition has a normally three-framed \(T^3\) inverse image.
The calculation above is independent of its class and supplies an
actual obstruction to null-homotopy of the original \(q_2\), even
when that other postcomposition is null-homotopic.  The exact
morphism carrying this obstruction is
\[
 (p_3)_*:[X,\mathcal A]_*\longrightarrow[X,S^4]_*,\qquad
 [q_2]\longmapsto[H]=\eta_4\eta_5.
\]
No conclusion about an arbitrary earlier quotient follows merely
from the diffeomorphism type of one inverse image; the specified
four-framing is what proves this value.

\begin{corollary}[The exact order-two quotient detected by the original cusp map]
\label{wb:h24:ret:q2-detected-quotient}
Let \(d:(S^6,*)\to(X,x_{\mathrm{base}})\) be the fixed based
diffeomorphism, and put
\[
\alpha=[q_2\circ d]\in\pi_6(\mathcal A,v_0).
\]
The original top-cell postcomposition induces a surjective homomorphism
\[
\varepsilon=(p_3)_*:\pi_6(\mathcal A,v_0)
  \longrightarrow\pi_6(S^4,*)\cong\mathbb Z/2,
\qquad \varepsilon(\alpha)=\eta_4\circ\eta_5.
\]
It gives the exact quotient isomorphism
\[
\pi_6(\mathcal A,v_0)/\ker\varepsilon
 \xrightarrow{\ [\beta]\mapsto\varepsilon(\beta)\ }
\mathbb Z/2.
\]
Every odd multiple of \(\alpha\) is nonzero. Its order is either
infinite or an even positive integer; this calculation does not determine
that order or assert a splitting of the displayed quotient.
\end{corollary}
\begin{proof}
Postcomposition distributes over the pinch sum on the source sphere,
so \((p_3)_*\) is a group homomorphism. Theorem~\ref{maps:src:B.3.1}
gives \((p_3)_*(\alpha)=[p_3q_2d]=\eta_4\eta_5\ne0\).
This is the nonzero element of the cyclic group of order two, proving
surjectivity. The first isomorphism theorem gives the stated quotient
with exactly the displayed map. For every integer \(n\),
\[
\varepsilon(n\alpha)=(n\bmod2)\eta_4\eta_5.
\]
An odd multiple therefore has nonzero image and cannot vanish.
If \(m\alpha=0\), applying \(\varepsilon\) forces \(m\)
to be even. These arguments prove all the assertions without inferring
the exact order of \(\alpha\).
\end{proof}


